\chapter{Güven Aralığı Tahmini ve Güven Aralıkları; Güven Düzeyi ve Doğru Yorumlama}
\label{ch:confidence-intervals}

\begin{prismbox}[PrismLight]{Öğrenme çıktıları}
Bu bölümün sonunda:
\begin{itemize}
  \item ortalama, oran ve grup farkı için güven aralığı kurabilecek,
  \item güven düzeyi, standart hata ve örneklem hacminin aralık genişliğine etkisini açıklayabilecek,
  \item güven aralığının uzun dönemli kapsama anlamını doğru yorumlayabilecek,
  \item hata payına göre örneklem hacmi planlayıp sonucu uygulamalı önemle değerlendirebileceksiniz.
\end{itemize}
\end{prismbox}

\begin{prerequisitebox}
Nokta tahmini, standart hata, örnekleme dağılımı, normal dağılım ve Student
$t$ kritik değeri kavramları hatırlanmalıdır.
\end{prerequisitebox}

Nokta tahmini ve tahmin edicinin belirsizliği için Bölüm~\ref{ch:point-estimation},
standart hata ve $t$ yaklaşımının temeli için Bölüm~\ref{ch:clt-se} kullanılmalıdır.

\section{Nokta tahmininden aralık tahminine}

\begin{definitionbox}
Nokta tahmini parametre için tek bir en iyi değeri verir. Güven aralığı ise
tahminin örnekleme değişkenliğini görünür kılan, belirli bir kapsama yöntemiyle
üretilmiş alt ve üst sınırlar kümesidir.
\end{definitionbox}
Güven aralığının genel yapısı
\[
\text{tahmin}\ \pm\ \text{kritik değer}\times\text{standart hata}
\]
biçimindedir.

\section{Ortalama için güven aralığı}
\label{sec:ci-mean}

Evren standart sapması bilinmediğinde ve koşullar uygun olduğunda
\[
\bar x\pm t_{1-\alpha/2,n-1}\frac{s}{\sqrt n}
\]
aralığı kullanılır. Örneğin $\bar x=40$, $s=8$, $n=64$ ve yaklaşık kritik
değer 2 ise hata payı $2(8/8)=2$ ve yaklaşık yüzde 95 güven aralığı
$[38,42]$ olur.

\subsection{Çözümlü örnek}

$n=25$ öğrencinin puan ortalaması $\bar x=72$, standart sapması $s=10$
olsun. $sd=24$ için yüzde 95 güven düzeyindeki kritik değer yaklaşık
$t_{0{,}975,24}=2{,}064$'tür. Standart hata 2, hata payı
$2{,}064\times2=4{,}128$ olur. Güven aralığı
\[
72\pm4{,}128=[67{,}87;76{,}13]
\]
olarak bulunur. Sonuç, yöntemin varsayımları altında evren ortalaması için
makul değerlerin yaklaşık 67,87 ile 76,13 arasında olduğunu gösterir.

\section{Güven düzeyinin doğru anlamı}
\label{sec:ci-interpretation}

Yüzde 95 güven, hesaplanan belirli aralığın parametreyi yüzde 95 olasılıkla
içerdiği anlamına gelmez. Klasik yaklaşımda parametre sabittir. Aynı örnekleme
ve aralık kurma işlemi çok kez tekrarlandığında kurulan aralıkların yaklaşık
yüzde 95'i gerçek parametreyi kapsar. Güven aralıklarının yalnız ikili bir
``anlamlı/anlamsız'' kararı yerine tahmin belirsizliğini göstermede
kullanılması özellikle önerilmektedir \citep{cumming2014,appelbaum2018}.

\begin{figure}[H]
\centering
\begin{tikzpicture}[x=0.9cm,y=0.52cm]
  \draw[PrismBlue,very thick] (5,0.25) -- (5,10.55);
  \node[above,font=\scriptsize,text=PrismBlue] at (5,10.55) {gerçek $\mu$};
  \foreach \y/\lo/\hi/\m in {
    1/3.7/5.8/4.75,2/4.1/6.0/5.05,3/3.9/5.5/4.70,4/4.4/6.3/5.35,
    5/3.5/4.8/4.15,6/4.3/5.7/5.00,7/3.8/5.3/4.55,8/4.5/6.1/5.30,
    9/3.6/5.4/4.50,10/5.2/6.7/5.95}
    {\draw[PrismBlue,line width=1.2pt] (\lo,\y)--(\hi,\y); \fill[PrismBlue] (\m,\y) circle (1.7pt);}
  \draw[PrismWarnOrange!90!black,line width=1.6pt] (5.2,10)--(6.7,10);
  \fill[PrismWarnOrange!90!black] (5.95,10) circle (1.9pt);
  \node[right,font=\scriptsize,text=PrismWarnOrange!80!black] at (6.75,10) {kapsamıyor};
\end{tikzpicture}
\caption{Tekrarlı örneklemlerde güven aralıklarının kapsama mantığı.}
\label{fig:ci-coverage}
\end{figure}

Şekil~\ref{fig:ci-coverage}'de okunabilirlik için yalnız on tekrar gösterilmiştir.
Yüzde 95 kapsama
özelliği on aralıktan mutlaka 9 veya 10'unun değil, çok sayıda tekrar boyunca
aralıkların yaklaşık yüzde 95'inin parametreyi kapsamasını ifade eder.

\begin{mistakebox}
Güven aralığının $H_0$ değerini içermesi, iki grubun aynı olduğunun kanıtı
değildir. Veri, farkın sıfır olmasıyla uyumludur; fakat aralığın genişliği
önemli büyüklükteki farkları da içeriyor olabilir.
\end{mistakebox}

\section{Aralık genişliğini belirleyen etkenler}

\begin{itemize}
  \item Güven düzeyi yükseldikçe kritik değer ve aralık genişliği artar.
  \item Örneklem hacmi arttıkça standart hata ve aralık genişliği azalır.
  \item Verideki değişkenlik arttıkça aralık genişler.
  \item Kümeleme, ağırlıklandırma ve bağımlı gözlemler standart hata hesabını değiştirir.
\end{itemize}

\section{Oran ve fark için aralıklar}
\label{sec:ci-proportion-difference}

Örneklem oranı için büyük örneklem yaklaşımında standart hata
\[
\SE(\hat p)=\sqrt{\frac{\hat p(1-\hat p)}{n}}
\]
olur. İki grubun karşılaştırılmasında yalnız ayrı grup aralıklarına bakmak
yerine doğrudan $\mu_1-\mu_2$ veya $p_1-p_2$ farkı için güven aralığı
hesaplanmalıdır.

Bağımsız iki grubun ortalama farkı için yaklaşık standart hata
\[
\SE(\bar X_1-\bar X_2)
=\sqrt{\frac{s_1^2}{n_1}+\frac{s_2^2}{n_2}}
\]
biçimindedir. Eşleştirilmiş ölçümlerde ise iki ayrı grup ortalaması yerine
her birim için hesaplanan farkların ortalaması kullanılır:
\[
\bar d\pm t_{1-\alpha/2,n-1}\frac{s_d}{\sqrt n}.
\]
Bağımsız ve eşleştirilmiş tasarımlar aynı veri yapısı değildir; standart hata
araştırma tasarımına göre seçilmelidir.

\section{İstatistiksel ve uygulamalı önem}
\label{sec:ci-practical}

Aralığın sıfırı dışlaması istatistiksel kanıta, genişliği tahminin
hassasiyetine, uç noktaları ise uygulamada makul etkilerin aralığına ilişkin
bilgi verir. Sonuç, yalnız ``anlamlı/anlamsız'' biçiminde sınıflandırılmamalıdır.

\section{Planlama bakımından hata payı}

Ortalama için normal yaklaşım altında hedef hata payı $h$ ise yaklaşık
örneklem hacmi
\[
n\approx\left(\frac{z_{1-\alpha/2}\sigma}{h}\right)^2
\]
ile planlanabilir. Kullanılan $\sigma$ değeri önceki araştırma veya pilot
çalışmadan gelmelidir. Tasarım etkisi, beklenen yanıtlamama ve alt grup
analizleri varsa ham hesap artırılmalıdır. Örneklem planında bilimsel olarak
önemli en küçük etkinin önceden belirlenmesi için bkz. \citet{lakens2022}.

\begin{outputguidebox}
Yazılım çıktısında güven düzeyini, tahmini, alt ve üst sınırı, kullanılan
standart hata yöntemini ve örneklem hacmini kontrol edin. Aralık ters yönde
yazılmış bir fark tanımına dayanıyorsa işaretin yorumu da tersine döner.
\end{outputguidebox}

\section{Yazılım uygulaması: ortalama güven aralığı}

\begin{softwarebox}
\begin{lstlisting}
from scipy import stats
import numpy as np

n, mean, sd = 25, 72, 10
se = sd / np.sqrt(n)
interval = stats.t.interval(
    confidence=.95, df=n-1, loc=mean, scale=se)
print(interval)
\end{lstlisting}
Çıktı yaklaşık $(67{,}87;76{,}13)$ değerlerini verir. Kodda güven düzeyi,
serbestlik derecesi, merkez ve standart hata ayrı ayrı görülebilmelidir.
\end{softwarebox}

\begin{assumptionbox}
Ortalama için t aralığı gözlemlerin bağımsızlığını ve özellikle küçük
örneklemlerde dağılımın aşırı çarpık ya da etkili aykırı değerli olmamasını
gerektirir. Güven aralığı seçim yanlılığını veya ölçüm hatasını kapsamaz.
\end{assumptionbox}

\section{Çözümlü problemler}

\subsection*{Problem 1: Ortalama için t güven aralığı}

Bağımsız seçilen 36 öğrencinin bir ölçekten aldığı puanların ortalaması 85,
standart sapması 12'dir. $sd=35$ için yüzde 95 kritik değerini
$t_{0{,}975,35}=2{,}030$ alarak evren ortalaması için güven aralığını bulun.

\textbf{Çözüm.} Evren standart sapması bilinmediği için t aralığı kullanılır.
Standart hata
\[
\SE(\bar X)=\frac{12}{\sqrt{36}}=2
\]
ve hata payı
\[
h=2{,}030(2)=4{,}06
\]
olur. Güven aralığı
\[
85\pm4{,}06=[80{,}94;89{,}06]
\]
biçimindedir.

\textbf{Yorum.} Aynı örnekleme ve aralık kurma yöntemi çok kez
tekrarlansaydı, elde edilen aralıkların yaklaşık yüzde 95'i gerçek evren
ortalamasını kapsardı. Gözlenen veride evren ortalaması için makul değerler
yaklaşık 80,94 ile 89,06 arasındadır.

\subsection*{Problem 2: Oran için büyük örneklem güven aralığı}

Rastgele seçilen 400 öğrencinin 260'ı bir dijital öğrenme aracını yararlı
bulmuştur. Evren oranı için yaklaşık yüzde 95 güven aralığını hesaplayın.

\textbf{Çözüm.} Örneklem oranı
\[
\hat p=\frac{260}{400}=0{,}65
\]
ve tahmini standart hata
\[
\SE(\hat p)=\sqrt{\frac{0{,}65(0{,}35)}{400}}
\approx0{,}02385
\]
olur. Yüzde 95 normal kritik değeri 1,96 ile hata payı
\[
h=1{,}96(0{,}02385)\approx0{,}0467
\]
ve güven aralığı
\[
0{,}65\pm0{,}0467=[0{,}603;0{,}697]
\]
olarak bulunur. Yaklaşık yüzde 95 aralık yüzde 60,3 ile yüzde 69,7'dir.

Başarı ve başarısızlık sayıları sırasıyla 260 ve 140 olduğundan büyük örneklem
yaklaşımı için sayısal koşullar güçlüdür. Küçük örneklemde veya oran sıfıra ya
da bire yakınsa Wilson ya da kesin aralık gibi daha uygun yöntemler
düşünülmelidir.

\subsection*{Problem 3: İki bağımsız ortalama farkı}

A öğretim yöntemini kullanan grubun ortalaması 78, B grubunun ortalaması
72'dir. $\mu_A-\mu_B$ farkının tahmini standart hatası 2,5 ve uygun serbestlik
derecesi için yüzde 95 kritik t değeri 2,021 olsun. Fark için güven aralığını
bulup yorumlayın.

\textbf{Çözüm.} Gözlenen ortalama farkı
\[
\bar x_A-\bar x_B=78-72=6
\]
puandır. Hata payı
\[
h=2{,}021(2{,}5)=5{,}0525
\]
olduğundan güven aralığı
\[
6\pm5{,}0525=[0{,}95;11{,}05]
\]
biçimindedir. Aralık sıfırı dışladığı için aynı çift yönlü test yaklaşık
$\alpha=0{,}05$ düzeyinde fark için kanıt verir. Bununla birlikte gerçek
farkın bir puana yakın da 11 puana yakın da olabileceği görülmektedir;
uygulamalı önem bu aralığın tamamı dikkate alınarak değerlendirilmelidir.

Fark B eksi A olarak tanımlansaydı tahmin ve uç noktaların işareti değişir,
aralık $[-11{,}05;-0{,}95]$ olurdu; bilimsel sonuç değişmezdi.

\subsection*{Problem 4: Güven düzeyi ile aralık genişliği}

Bir tahmin 50 ve standart hatası 2 olsun. Normal yaklaşım altında yüzde 90,
yüzde 95 ve yüzde 99 güven aralıklarını sırasıyla 1,645, 1,96 ve 2,576 kritik
değerleriyle hesaplayın.

\textbf{Çözüm.}

\begin{center}
\small
\begin{tabular}{lrrl}
\toprule
Güven düzeyi & Kritik değer & Hata payı & Güven aralığı \\
\midrule
Yüzde 90 & 1,645 & 3,29 & $[46{,}71;53{,}29]$ \\
Yüzde 95 & 1,960 & 3,92 & $[46{,}08;53{,}92]$ \\
Yüzde 99 & 2,576 & 5,15 & $[44{,}85;55{,}15]$ \\
\bottomrule
\end{tabular}
\end{center}

Standart hata ve tahmin değişmediği hâlde daha yüksek güven düzeyi daha büyük
kritik değer ve daha geniş aralık üretmiştir. Daha yüksek kapsama hedefinin
maliyeti daha düşük hassasiyettir.

\subsection*{Problem 5: Hedef hata payına göre örneklem planı}

Bir başarı ölçeğinin standart sapması önceki çalışmalarda yaklaşık 12 puan
bulunmuştur. Yüzde 95 güven düzeyinde ortalama için hata payının en fazla 2
puan olması isteniyor. Gerekli örneklem hacmini bulun. Yüzde 20 yanıtlamama
bekleniyorsa kaç kişi çalışmaya davet edilmelidir?

\textbf{Çözüm.} Normal yaklaşım altında
\[
n\approx\left(\frac{1{,}96(12)}{2}\right)^2
=11{,}76^2=138{,}30
\]
olur. Hedef hata payını aşmamak için yukarı yuvarlanır ve en az 139 geçerli
gözlem gerekir.

Yanıt oranının yüzde 80 olması bekleniyorsa davet edilecek sayı
\[
n_{\mathrm{davet}}=\frac{139}{0{,}80}=173{,}75
\]
olduğundan en az 174 kişiye ulaşılması planlanmalıdır. Bu artış yalnız beklenen
yanıt kaybını karşılar; küme tasarımı veya alt grup analizleri varsa ek
düzeltme gerekir.

\subsection*{Problem 6: Sıfırı içeren aralığın yorumu}

İki öğretim yönteminin ortalama puan farkı A eksi B olarak 3,5 puan, yüzde 95
güven aralığı $[-1{,}5;8{,}5]$ bulunmuştur. Uygulamada en az 5 puanlık farkın
önemli kabul edildiği varsayılsın. Sonucu istatistiksel ve uygulamalı açıdan
yorumlayın.

\textbf{Çözüm.} Aralık sıfırı içerdiği için yüzde 5 düzeyindeki çift yönlü
testte sıfır fark hipotezi reddedilemez. Bu sonuç yöntemlerin eşit olduğunu
kanıtlamaz. Aralık B'nin A'dan 1,5 puan daha iyi olmasından A'nın B'den 8,5
puan daha iyi olmasına kadar geniş bir değer kümesini içermektedir.

Uygulamalı önem eşiği 5 puan olduğundan aralık hem önemsiz farkları hem A
lehine önemli kabul edilecek farkları kapsamaktadır. Dolayısıyla çalışma
``etki yok'' sonucundan çok, yön ve büyüklük hakkında yetersiz hassasiyet
göstermektedir. Daha büyük veya daha iyi tasarlanmış bir çalışma belirsizliği
azaltabilir.

\begin{mistakebox}
Sıfırı içeren geniş bir güven aralığını ``iki yöntem aynıdır'' biçiminde
yorumlamak hatalıdır. Eşdeğerlik göstermek için önceden belirlenmiş uygulamalı
sınırlar ve buna uygun eşdeğerlik analizi gerekir.
\end{mistakebox}

\begin{outputguidebox}
Güven aralığı çıktısında önce farkın hangi yönde tanımlandığını, güven
düzeyini ve kullanılan standart hata yöntemini kontrol edin. Ardından aralığın
sıfır gibi referans değerleri ve uygulamalı önem sınırlarını içerip
içermediğini değerlendirin. Yalnız \pval-değerine göre yorum yapmayın.
\end{outputguidebox}

\section{Literatür verisiyle yazılım uygulamaları}
\label{sec:spss-b08}

\textbf{Amaç ve veri.} \texttt{b08.xlsx} ile ortalama için $t$ güven
aralığını, \texttt{pass10} oranı için Wilson aralığını bulun
\citep{cortez2008data}.

\textbf{Ortak veri.} Üç yazılımda aynı kayıtlar ve değişkenler kullanılır. R ve Python için \texttt{companion/spss/csv/b08.csv}, Excel dosyasının aynı veri sayfasını içerir. Kodları proje kökünden çalıştırın; veri sözlüğü ve hazırlık için bkz. sayfa~\pageref{sec:spss-guide}.

\subsection{IBM SPSS uygulaması}
\textbf{Başlamadan önce.} Rehberdeki veri açma adımıyla yalnız \texttt{open-b08.sps} dosyasını çalıştırın; etiketli veri açılır. Menü yolu ve aşağıdaki tam komut dosyası alternatif uygulama yollarıdır.

\begin{spssadimlar}
\spssadim{\emph{Analyze $\to$ Descriptive Statistics $\to$ Explore} yolunda \texttt{g3} değişkenini \emph{Dependent List} alanına taşıyın. \emph{Factor List} boş kalsın.}
\spssadim{\emph{Display} için \emph{Statistics} seçin. \emph{Statistics} penceresinde \emph{Descriptives} ve yüzde 95 ortalama güven düzeyini belirleyin; \emph{Continue $\to$ OK} seçin.}
\spssadim{Çıktıdaki \emph{95\% Confidence Interval for Mean} altında \emph{Lower Bound} ve \emph{Upper Bound} satırlarını bulun. Bunlar not ortalamasının sınırlarıdır.}
\spssadim{\emph{Analyze $\to$ Descriptive Statistics $\to$ Frequencies} yolunda \texttt{pass10} tablosunu alın; 1 kategorisinin frekansını ve toplamı okuyun.}
\spssadim{Oran için Wilson aralığı, Explore içindeki ortalama aralığı değildir. Aşağıdaki tam komut dosyasını \emph{File $\to$ Open $\to$ Syntax} ile açıp \emph{Run $\to$ All} seçin; \texttt{alt} ve \texttt{ust} sınırlarını \texttt{LIST} çıktısından okuyun.}
\end{spssadimlar}

{\raggedright
\noindent\textbf{Komutların görevi.} \texttt{EXAMINE} ortalamanın $t$ aralığını verir. Sonraki \texttt{AGGREGATE} ve \texttt{COMPUTE} komutları oran için Wilson sınırlarını hesaplar; bunlar ayrı çıktılardır.\par}

\par\noindent\begin{minipage}{\linewidth}
\textbf{Uygulama kodu:} \texttt{b08}. Veri: \texttt{excel/b08.xlsx};
komut: \texttt{syntax/b08.sps} (\texttt{companion/spss} altında).

% Derleme harici .sps dosyasına bağlı değildir. Eşlikçi kopya:
% companion/spss/syntax/b08.sps
\begin{lstlisting}[style=prismSPSS]
INSERT FILE='companion/spss/syntax/open-b08.sps'
 ERROR=STOP.
EXAMINE VARIABLES=g3
 /PLOT=NONE /STATISTICS=DESCRIPTIVES
 /CINTERVAL=95.
FREQUENCIES VARIABLES=pass10.
AGGREGATE /OUTFILE=*
 /BREAK= /n=N /oran=MEAN(pass10).
COMPUTE z=1.959963984540054.
COMPUTE payda=1+z**2/n.
COMPUTE merkez=(oran+z**2/(2*n))/payda.
COMPUTE yari=z*SQRT(oran*(1-oran)/n+z**2/(4*n**2))/payda.
COMPUTE alt=merkez-yari.
COMPUTE ust=merkez+yari.
FORMATS oran alt ust (F10.4).
LIST VARIABLES=n oran alt ust.
\end{lstlisting}

\end{minipage}\par

\begin{outputguidebox}
\textbf{Beklenen kontrol değerleri (SPSS burada çalıştırılmadı).}
\emph{Descriptives} tablosunda ortalamanın yüzde 95 güven aralığı
$[\SPMatCILow;\SPMatCIHigh]$ olur. \texttt{LIST} çıktısında oran
\SPMatPHat, Wilson alt ve üst sınırları \SPWilsonLow\ ve \SPWilsonHigh'dır.
Birinci aralığın birimi not puanı, ikincisininki orandır.
Bu aralıklar öğrencilerin yüzde 95'inin notlarını kapsayan aralıklar değildir.
Model ve örnekleme varsayımları geçerliyken tekrarlanan örneklemelerden elde
edilen aralıkların kapsama özelliğiyle yorumlanırlar. İki okuldan veri
toplanmış olmasının genelleme sınırlılığı aralık hesaplayarak ortadan kalkmaz.
\end{outputguidebox}


\par\noindent\begin{minipage}{\linewidth}
\subsection{R uygulaması}
\label{sec:r-gercek-b08}

İlk satır yalnız ortalama için $t$ aralığını seçer. Sonraki hesap, SPSS ile aynı Wilson formülünü kullanır.

\begin{lstlisting}[style=prismR]
d <- read.csv("companion/spss/csv/b08.csv")
stopifnot(!anyNA(d))
print(t.test(d$g3, conf.level=.95)$conf.int)
print(table(d$pass10))
n <- nrow(d); p <- mean(d$pass10)
z <- qnorm(.975); den <- 1 + z^2/n
center <- (p + z^2/(2*n))/den
half <- z * sqrt(p*(1-p)/n + z^2/(4*n^2))/den
print(c(proportion=p, lower=center-half,
        upper=center+half))
\end{lstlisting}

\textbf{Çıktıyı okuma.} Aşağıdaki ortak karşılaştırmada verilen nicelikleri, kodun yazdırdığı etiketler ve değişken sırasıyla eşleştirin. R kodu bu ortamda çalıştırılmamıştır.
\end{minipage}\par

\par\noindent\begin{minipage}{\linewidth}
\subsection{Python uygulaması}
\label{sec:python-b08}

Ortalama aralığı ile Wilson oran aralığı ayrı hesaplanır. \texttt{wilson} alt ve üst sınırı saklar.

\begin{lstlisting}[style=prismPythonApp]
import pandas as pd
import numpy as np
from scipy import stats

d = pd.read_csv("companion/spss/csv/b08.csv")
assert not d.isna().any().any()
mean_ci = stats.ttest_1samp(d.g3, 0).confidence_interval(.95)
print(mean_ci)
print(d.pass10.value_counts().sort_index())
n, p = len(d), d.pass10.mean()
z = stats.norm.ppf(.975); den = 1 + z*z/n
center = (p + z*z/(2*n))/den
half = z*np.sqrt(p*(1-p)/n + z*z/(4*n*n))/den
wilson = (center-half, center+half)
print(p, wilson)
\end{lstlisting}

\textbf{Çıktıyı okuma.} Yazdırılan değerleri aşağıdaki ortak kontrol noktalarıyla karşılaştırın; yalnız testin $p$ değerine bakarak rapor yazmayın.
\end{minipage}\par

\subsection{Sonuçların karşılaştırılması ve ortak rapor}
\begin{outputguidebox}
Ortalama aralığı $[\SPMatCILow;\SPMatCIHigh]$, Wilson aralığı $[\SPWilsonLow;\SPWilsonHigh]$ olmalıdır. Wald veya süreklilik düzeltmeli bir aralıkla karşılaştırma yapılmamalıdır.
\end{outputguidebox}

\textbf{Raporlama örneği (kontrol değerleriyle).}
``Bağımsız gözlemler modeli altında ortalama için yüzde 95 güven aralığı $[\SPMatCILow;\SPMatCIHigh]$, eşik göstergesinin oranı için Wilson aralığı $[\SPWilsonLow;\SPWilsonHigh]$ bulunmuştur.''


\textbf{Görev.} Ortalama aralığının 10'u içerip içermediğini bulun.
Aynı veri, çift yönlü hipotez ve anlamlılık düzeyinde bunun $t$ testiyle
ilişkisini açıklayın. \texttt{.95} güven olasılığı ile \texttt{95} yüzde
gösteriminin farklı komutlarda kullanılmasına dikkat edin.

\section{R ile çözümlü uygulama}
\label{sec:r-b08}

\textbf{Soru.} $n=25$, $\bar x=72$ ve $s=10$ için evren ortalamasının
yüzde 95 ve yüzde 99 güven aralıklarını hesaplayın. Burada ham gözlemler
değil özet istatistikler verildiği için $t$ kritik değeriyle doğrudan hesap yapın.

\begin{lstlisting}[style=prismR]
hacim <- 25
ortalama <- 72
standart_sapma <- 10
standart_hata <- standart_sapma / sqrt(hacim)
serbestlik <- hacim - 1
guven_araligi <- function(duzey) {
  kritik <- qt((1 + duzey) / 2, df = serbestlik)
  ortalama + c(-1, 1) * kritik * standart_hata
}
c(standart_hata = standart_hata,
  kritik_95 = qt(0.975, df = serbestlik))
rbind(GA95 = guven_araligi(0.95),
      GA99 = guven_araligi(0.99))
\end{lstlisting}

\begin{outputguidebox}
\textbf{Çözüm ve kontrol değerleri.} Standart hata $10/\sqrt{25}=2$,
serbestlik derecesi 24 ve yüzde 95 aralık için kritik değer 2,0639'dur.
Yüzde 95 güven aralığı $[67{,}8722;76{,}1278]$, yüzde 99 güven aralığı
yaklaşık $[66{,}4061;77{,}5939]$ olur.
\end{outputguidebox}

\textbf{Yorum.} Bağımsız gözlemler ve küçük örneklemde yaklaşık normal
bir evren modeli altında bu $t$ aralığı kullanılır. Özet istatistikler aykırı
değerleri ve dağılım biçimini denetlemeye yetmez. Ham veri elde olsaydı
\texttt{t.test(puan, conf.level = 0.95)\$conf.int} kullanılabilirdi;
\texttt{puan} yerine üç özet sayıyı koymak yanlış olurdu.

\textbf{Pekiştirme ve yanıt.} Güven düzeyi arttıkça kritik değer ve aralık
genişliği artar. Yüzde 95 ifadesi, sabit parametrenin hesaplanmış bu aralıkta
bulunma olasılığı değil yöntemin uzun dönem kapsama oranıdır.

\section{Bölüm özeti}

\begin{itemize}
  \item Güven aralığı tahmin, kritik değer ve standart hatayı birleştirir.
  \item Güven düzeyi belirli tek bir aralığa değil, aralık kurma yönteminin uzun dönem başarısına ilişkindir.
  \item Güven düzeyi, değişkenlik ve örneklem hacmi aralık genişliğini belirler.
  \item Aralık hem istatistiksel kanıtı hem uygulamada makul etki büyüklüklerini gösterir.
  \item Grup karşılaştırmalarında ayrı grup aralıkları yerine doğrudan fark
  için güven aralığı kurulmalıdır.
  \item Sıfırı içeren geniş bir aralık eşitliği kanıtlamaz; çoğu zaman
  çalışmanın hassasiyetinin sınırlı olduğunu gösterir.
\end{itemize}

\section*{Alıştırmalar}

\begin{enumerate}
  \item Güven düzeyi yüzde 95'ten yüzde 99'a çıkarılırsa aralık nasıl değişir?
  \item Örneklem hacmi dört katına çıkarsa hata payı yaklaşık nasıl değişir?
  \item Güven aralığının hem istatistiksel hem uygulamalı yorumunu açıklayın.
  \item Çözümlü örnekte standart sapma 15 olsaydı aralık nasıl değişirdi?
  \item Güven aralığının dar olmasına rağmen neden yanıltıcı olabileceği bir durum yazın.
  \item $n=49$, $\bar x=60$, $s=14$ ve kritik değer 2,01 için ortalama güven
  aralığını hesaplayın.
  \item 500 öğrencinin 310'unun olumlu yanıt verdiği bir araştırmada oran için
  yaklaşık yüzde 95 güven aralığını bulun.
  \item Tahmin 24 ve standart hata 3 ise yüzde 90 ve yüzde 95 normal güven
  aralıklarını karşılaştırın.
  \item Ortalama farkı 4, standart hatası 1,5 ve kritik t değeri 2 ise fark
  güven aralığını hesaplayıp sıfır bakımından yorumlayın.
  \item $\sigma=16$, hedef hata payı 2 ve yüzde 95 güven düzeyi için yaklaşık
  örneklem hacmini belirleyin.
  \item Gerekli geçerli örneklem 240 ve beklenen yanıt oranı yüzde 75 ise kaç
  kişinin davet edilmesi gerektiğini hesaplayın.
  \item Yüzde 99 güven aralığının neden yüzde 95 aralığından daha geniş
  olduğunu kapsama--hassasiyet dengesiyle açıklayın.
  \item $[-2;12]$ biçimindeki bir fark aralığını istatistiksel anlamlılık,
  etki yönü ve uygulamalı önem bakımından yorumlayın.
  \item Ayrı iki grup güven aralığının örtüşmesine bakmanın neden doğrudan fark
  aralığının yerini tutmadığını açıklayın.
\end{enumerate}